\scanpage{318}
\mspage{268}

ce qui importe pour la validité de cet énoncé, c'est (on s'en doute) que $Z_{\unsure{*|\bullet}}$ soit $W$-cocartésien. Donc il faut d'abord prouver ceci

\uline{Lemme} Si le \deleted{\unsure{cub}}diagramme cubique p. 265 est tel que les faces $X,Y$ soient $W$-cocartésiens, et \deleted{\unsure{$f_i$}} les $f_i$ ($i\in\{0,1,2,3\}$) des $W$-cofibrations, alors le carré $Z$ \deleted{qui s'en \unsure{déduit} en prenant} des cofibres des $f_i$ est lui aussi $W$-cocartésien.

Ce lemme résulte de la théorie du ``déploiement des diagrammes cubiques'', \unsure{relatif} \deleted{\unsure{une}} à une catégorie de modèles $(\mathcal M,W)$ admettant une théorie des cofibres $W$-homotopiques, qu'il me faudra bien écrire ! (Mais l'existence \inserted{\deleted{\unsure{pour}} $(\mathrm{Cat},W)$} d'une telle théorie des cofibres $W$-homotopiques, et par suite de celle du déploiement des diagrammes dans $\mathcal M$ de type $\unsure{\Delta_1^{I}|\Delta^{1\,I}}$, $I$ ens. d'indices \unsure{finis|fini}, ne dépend que de $W(7\ \mathrm{bis})$ -- on n'a pas \unsure{à} faire appel à $W(7)$.)
% archive sheet number: 318
